\chapter{Introduction}\label{chap:intro}
\section{Context of the Thesis}\label{sec:context}

Machine Learning (ML) has become a central component of many digital services and supports decisions and predictions in many domains. At the core of these capabilities lies the ability of learning algorithms to extract statistical properties from data and use them to generalize beyond the observations available during training. In many learning settings, larger and more diverse datasets allow models to capture richer patterns and improve their predictive capabilities~\cite{banko-brill-2001-scaling,shalev2014understanding,hestness2017deeplearningscalingpredictable,kaplan2020scalinglawsneurallanguage}. The growing availability of digital data has therefore become one of the main resources enabling the development of increasingly capable machine-learning systems.

Most of the information used to train these models originates from human activity and describes individuals, either directly or indirectly. A sequence of locations may reveal where someone lives or works; consumption and browsing histories may expose interests, habits, or preferences; medical and financial records describe particularly sensitive aspects of a person's life.
Even if harmless alone, when combinaed, these infotmation traces can become sensitive~\cite{narayanan2008robust}. From the perspective of a learning algorithm, individual records form a dataset from which population-level patterns can be inferred, but, from the perspective of the individual, each observation may correspond to a fragment of their private life.


For this reason, personal information is subject to legal requirements governing its collection and use, including principles such as purpose limitation and data minimization, as well as rights granted to the individuals concerned~\cite{eu_gdpr_2016,california_ccpa_cpra}. More recent regulatory frameworks further address the use of sensitive data and the deployment of data-driven and machine-learning systems in specific domains~\cite{eu_ai_act_2024,eu_ehds_2025,Buick_copyright}. Yet legal safeguards do not eliminate the underlying imbalance between the individuals who generate the data and the organizations that use them to provide digital services. Once data leave the direct control of their owners, those individuals lose control on how the information may later be retained, combined with other sources, or exploited.

This imbalance is particularly significant in modern digital ecosystems, where a relatively small number of organizations operate many of the devices, applications, platforms, and cloud services through which everyday activities are mediated. Their position gives them access to detailed information about users' behavior, preferences, movements, purchases, communications, and social interactions. Such knowledge enables valuable and highly personalized services, but can also be used for profiling, targeting, ranking, persuasion, facial recognition, and other forms of automated decision-making that may remain opaque to the people affected~\cite{pasquale_black_box_society_2015,veliz_privacy_is_power_2021}. Privacy therefore becomes not only a question of what information is disclosed, but also of who controls it, what can be inferred from it, and how much trust must be placed in the entities that hold it.


 Recent enforcement actions involving major technology companies demonstrates that privacy risks often arise not only from external attacks or accidental data breaches, but also from ordinary platform practices.
The Cambridge Analytica scandal provides a prominent example, where personal information collected through a digital platform was subsequently exploited for political profiling and targeted campaigning, raising concerns about consent, transparency, and secondary uses of personal data~\cite{ico_democracy_disrupted_2018,ico_data_analytics_political_campaigns_2018}. Similar concerns have emerged from enforcement actions involving major technology companies, including cases related to advertising practices, children's data, voice recordings, and the reuse of information collected for unauthorized purposes~\cite{cnil_google_gdpr_2019,ftc_youtube_coppa_2019,ftc_twitter_security_data_ads_2022,ftc_amazon_alexa_coppa_2023,texas_ag_google_privacy_2025}. 
Similar concerns arise in machine learning, where the demand for large and diverse training datasets creates strong incentives to collect and reuse data at scale. Controversies involving biometric systems trained on scraped or improperly obtained facial images \citeweb{noyb_clearview_criminal_complaint_2025}, copyrighted material used for model training \citeweb{reisner_ai_pirated_books_problem_2025}, and wearable devices capturing data from both users and bystanders \citeweb{svd_meta_ai_glasses_privacy_2026} demonstrate how the demand for data can blur the boundary between legitimate use and unauthorized collection or reuse.


Similar concerns arise in machine learning, where the demand for large and diverse training datasets creates strong incentives to collect and reuse data at scale. Controversies involving biometric systems trained on scraped or improperly obtained facial images~\citeweb{noyb_clearview_criminal_complaint_2025}, copyrighted material used for model training~\citeweb{reisner_ai_pirated_books_problem_2025}, and wearable devices capturing data from both users and bystanders~\citeweb{svd_meta_ai_glasses_privacy_2026} illustrate how the demand for data can blur the boundary between legitimate use and unauthorized collection. More generally, conventional centralized machine learning requires data held by different individuals or organizations to be brought together before they can be jointly processed. Even when such collection is legitimate, concentrating raw information under the control of a single entity reinforces the dependence on that entity to handle it appropriately.

A natural way to reduce this dependence is to limit the amount of raw information that must be disclosed in the first place. Rather than moving data towards a central learning infrastructure, the computation can instead be moved towards the data, allowing different parties to collaborate while retaining their records locally. The principle is particularly relevant when data are already distributed across personal devices, organizations, or institutions and cannot, or should not, be collected. Federated Learning (FL) realizes this approach~\cite{pmlr-v54-mcmahan17a}, enabling multiple participants to collaboratively train a shared model without requiring their raw datasets to be collected in a central repository.

A natural way to reduce this dependence is to limit the amount of raw information that must be disclosed in the first place. Rather than moving data towards a central learning infrastructure, computation can be moved closer to where the data are generated or stored, allowing different parties to collaborate without pooling their records. The principle is particularly relevant when information is already distributed across personal devices, organizations, or institutions and cannot, or should not, be centrally collected. Federated Learning (FL) provides a prominent framework for learning under these constraints~\cite{pmlr-v54-mcmahan17a}.


\section{Federated Learning} \label{sec:fed_learning}

Federated Learning enables multiple participants to train a shared machine-learning model while keeping their datasets at their original locations. Each client trains the current model locally on its own data and shares only the resulting model update, such as a gradient or updated model parameters~\cite{konecny2015federatedoptimizationdistributedoptimizationdatacenter, pmlr-v54-mcmahan17a}.
The main advantage and motivation behind  FL is its its natural adaptation to large-scale data collection settings. In applications involving mobile and edge devices, information is naturally produced across a large population of geographically distributed clients, while transferring and persistently storing all of it in a central infrastructure may be costly or impractical. 

 In the canonical FL formulation, a central server orchestrates the training process: it distributes the current model to selected clients, receives their locally computed updates, and aggregates them to obtain the next global model. A fully decentralized alternative removes the central coordinator and lets participants exchange information directly according to a communication graph, collectively updating their models through local interactions. Both settings avoid centralizing the underlying datasets, but differ in who observes the exchanged information and in the trust relationships created by the learning protocol. Throughout this thesis, we focus on the server-coordinated formulation of FL~\cite{pmlr-v54-mcmahan17a}.

A second important aspect of the server-coordinated approach is that it  limits the direct disclosure of personal information. Keeping records at their original location is naturally aligned with the principle of \emph{data minimization}, according to which personal information should be limited to what is necessary for a specified purpose~\cite{eu_gdpr_2016,eu_regulation_2018_1725,TRUONG2021102402}. This characteristic has made FL particularly attractive in privacy-sensitive domains such as healthcare and finance~\cite{Rieke2020,ZHANG2021106775,10.1145/3298981}. The privacy intuition is also reflected in the original formulation of FL, where the authors observe that the updates exchanged during training ``will generally contain much less information than the raw training data'' \cite{pmlr-v54-mcmahan17a}. Yet reducing direct access to the records is not equivalent to guaranteeing their privacy: the information exchanged during training remains dependent on the local data. We return to this distinction in Section~\ref{sec:privacy_in_fl}.




